\section{Background}\label{sec:background}

\subsection{Reinforcement Learning and Soft Actor-Critic}
Reinforcement learning (RL) models sequential decision-making as a Markov Decision Process
(MDP): at each step the agent observes a state, takes an action, and receives a reward, and it
learns a policy that maximises the expected discounted sum of future rewards
\cite{sutton2018rl}. In mobile-robot navigation, the state usually combines a depth image with
the position of the goal relative to the robot, and the action is a continuous linear and
angular velocity \cite{chaffre2020sim}.

Soft Actor-Critic (SAC) \cite{haarnoja2018sac} is the algorithm used in this study. Besides
maximising reward, it rewards the policy for staying random by adding an entropy term to the
objective:
\begin{equation}
J(\pi)=\sum_{t=0}^{T}\mathbb{E}_{(s_t,a_t)\sim\rho_\pi}\big[r(s_t,a_t)+\alpha\,\mathcal{H}(\pi(\cdot\mid s_t))\big],
\label{eq:sac}
\end{equation}
where $\mathcal{H}$ is the entropy of the policy and the temperature $\alpha$ sets how much
exploration is valued relative to reward. A policy that does not commit too early to one way of
solving the task is less likely to overfit to the simulator. SAC is off-policy: it stores past
experience in a replay buffer and reuses it many times, learns two critics with slowly updated
target copies to avoid overestimating values, and tunes $\alpha$ automatically
\cite{haarnoja2018sac}. It suits this study because Pepper's actions are continuous, because
reusing experience matters when every simulated step runs in real time, and because it has
already been used for sim-to-real navigation \cite{chaffre2020sim}.

\subsection{The Sim-to-Real Gap and Domain Randomisation}
A policy trained in simulation often fails on the real robot, because the simulator is only an
approximation of reality \cite{zhao2020sim, peng2018dynamics}. Zhao et al.\ \cite{zhao2020sim}
divide this \emph{reality gap} into mismatches in sensing, such as camera behaviour, and in
actuation and dynamics, such as friction, mass and latency. This study uses \emph{zero-shot
transfer}: the policy runs on the real robot without any further training.

Domain randomisation varies the simulator's parameters during training, so that the range of
conditions the policy sees is wide enough to include the real world \cite{tobin2017domain}.
\emph{Dynamics randomisation} varies physical parameters such as friction and actuator delay
\cite{peng2018dynamics}, and \emph{observation randomisation} varies what the sensors report,
for example by adding noise or dropping readings. Too much randomisation can harm rather than
help real-world performance \cite{zhao2020sim}, so this study sets the ranges from measurements
of the real robot where possible (Section~\ref{sec:methodology}).

\subsection{Photogrammetry}
Photogrammetry reconstructs a textured 3D model of a real space from many overlapping
photographs, by matching the same visual features across images taken from different viewpoints
and working out their 3D positions \cite{alhamadah2024photogrammetry}. The resulting mesh can be
imported into a robot simulator. Alhamadah et al.\ \cite{alhamadah2024photogrammetry} show that
consumer devices such as smartphones can produce such models with a mean dimensional error of
about 5\%.

\subsection{The Pepper Robot and ROS~2}
Pepper is a wheeled humanoid service robot made by SoftBank Robotics and widely used in research
\cite{pandey2018pepper, sekkat2024pepper}. Its hardware is too fragile for repeated real-world
training, so simulation is needed \cite{sekkat2024pepper}. Its depth camera has a narrow field
of view, and its visual localisation drifts over distance and is sensitive to lighting
\cite{gomez2018visual}, which is why this study uses depth images and a single room. The robot is
connected to the policy through ROS~2 \cite{macenski2022ros2}, whose named data channels
(topics) are used in the same way in the Gazebo simulator \cite{koenig2004gazebo} and on the real
Pepper, so the policy's inputs and outputs stay identical across the transfer.
